% ============================================================
%  Преамбула — заједничка подешавања за цео рад
%  Компајлира се искључиво са XeLaTeX (због fontspec/unicode-math)
% ============================================================

% ── језик и фонтови ─────────────────────────────────────────
\usepackage{fontspec}
\usepackage{polyglossia}
\setmainlanguage{serbian}
\setmainfont{Times New Roman}

% Моноспејс фонт мора да садржи ћирилицу (подразумевани Latin Modern Mono
% је нема, што polyglossia пријављује као грешку код \texttt{...}).
\setmonofont{Consolas}[Scale=0.92]

% ── математика ──────────────────────────────────────────────
% ВАЖНО: amsmath иде ПРЕ unicode-math, а amssymb се НЕ учитава --
% unicode-math већ дефинише све те симболе, па би amssymb изазвао
% сукоб („Command \eth already defined“).
\usepackage{amsmath}
\usepackage{unicode-math}
\setmathfont{Cambria Math}
\allowdisplaybreaks          % дозволи прелом вишередних формула преко стране

% ── распоред стране ─────────────────────────────────────────
% Маргине по шаблону ФОН-а: горња, доња и десна 2,5 цм; лева 3,1 цм.
\usepackage[a4paper,left=3.1cm,right=2.5cm,top=2.5cm,bottom=2.5cm]{geometry}
\usepackage{setspace}
\onehalfspacing               % проред између редова 1,5

% Проред параграфа по шаблону ФОН-а: 12pt пре и 6pt након параграфа.
% Word сажима суседне размаке на ВЕЋИ од та два (12pt), не на њихов
% збир, па је 12pt тачан еквивалент. LaTeX има само један размак
% између параграфа (\parskip), па се он поставља на 12pt.
% Растезање (plus/minus) намерно је нула: без тога LaTeX додатно
% развлачи размаке да попуни страну, што ствара неуједначене празнине.
\setlength{\parindent}{0pt}
\setlength{\parskip}{12pt}

% Обострано поравнање је подразумевано у LaTeX-у; овде се само
% ублажава превише агресивно растезање речи.
\usepackage{microtype}

% ── уједначавање вертикалних размака ────────────────────────
% Пошто је \parskip постављен на 12pt, LaTeX-ови подразумевани
% размаци око набрајања, формула и наслова СЕ ДОДАЈУ на њега и
% стварају видљиве празне редове. Зато се сви експлицитно своде
% на исти размак од 12pt, тако да документ има један једини
% вертикални ритам.

% Набрајања (itemize, enumerate, description)
\usepackage{enumitem}
\setlist{
  topsep    = 0pt,     % изнад/испод целе листе — наслеђује \parskip
  partopsep = 0pt,
  parsep    = 0pt,
  itemsep   = 6pt,     % између ставки: пола размака параграфа
  listparindent = 0pt
}

% Формуле: amsmath подразумевано додаје сопствени размак изнад и
% испод, који се сабира са \parskip. Своди се на исту вредност.
% Поставља се кроз \AtBeginDocument јер amsmath своје вредности
% дефинише при учитавању величине фонта и прегазио би ранију поставку.
% Формула је наставак реченице која јој претходи, а не нов параграф,
% па добија 6pt — пола размака параграфа. Са 12pt би изгледала као
% да је одвојена празним редом, што је управо оно што се избегава.
\AtBeginDocument{%
  \setlength{\abovedisplayskip}{6pt}%
  \setlength{\belowdisplayskip}{6pt}%
  \setlength{\abovedisplayshortskip}{6pt}%
  \setlength{\belowdisplayshortskip}{6pt}%
}

% Плутајући објекти (табеле): размак изнад и испод натписа и
% између објекта и текста.
\setlength{\textfloatsep}{12pt}
\setlength{\floatsep}{12pt}
\setlength{\intextsep}{12pt}
\setlength{\abovecaptionskip}{6pt}
\setlength{\belowcaptionskip}{6pt}

% ── табеле ──────────────────────────────────────────────────
\usepackage{booktabs}
\usepackage{tabularx}
\usepackage{longtable}
\usepackage{array}
\newcolumntype{L}[1]{>{\raggedright\arraybackslash}p{#1}}

% ── наслови по шаблону ФОН-а ────────────────────────────────
% Ниво 1: СВИМ ВЕЛИКИМ СЛОВИМА, фонт 14, затамњено.
% Ниво 2: СВИМ ВЕЛИКИМ СЛОВИМА, фонт 14 (без задебљања).
% Ниво 3: малим словима, фонт 14.
% Поглавља се не преламају на нову страну (шаблон нема преломе),
% па се користи \chapter* логика преко titlesec-а без \clearpage.
\usepackage{titlesec}

% Размак ИСПОД наслова постављен је на 0pt, а не на 6pt: \parskip од
% 12pt већ увлачи размак пре првог параграфа иза наслова, па би 6pt
% на то дало 18pt и видљиву празнину. Размак ИЗНАД наслова је 6pt,
% што уз \parskip претходног параграфа даје укупно 18pt раздвајања
% између претходног текста и новог наслова.
\titleformat{\chapter}[hang]
  {\fontsize{14}{17}\selectfont\bfseries\MakeUppercase}
  {\thechapter}{1em}{}
\titlespacing*{\chapter}{0pt}{0pt}{0pt}

\titleformat{\section}[hang]
  {\fontsize{14}{17}\selectfont\MakeUppercase}
  {\thesection}{1em}{}
\titlespacing*{\section}{0pt}{6pt}{0pt}

\titleformat{\subsection}[hang]
  {\fontsize{14}{17}\selectfont}
  {\thesubsection}{1em}{}
\titlespacing*{\subsection}{0pt}{6pt}{0pt}

% \paragraph се користи за истакнуте поднаслове у тексту (нпр. „Х1 --- ...“).
% Мора остати у истом реду са текстом који следи, без додатног размака.
\titleformat{\paragraph}[runin]
  {\normalfont\bfseries}{}{0pt}{}
\titlespacing*{\paragraph}{0pt}{0pt}{1em}

% Поглавље почиње на новој страни, али без празне стране између
% (oneside документ то и иначе поштује).
\let\originalchapter\chapter

% ── референце и хипервезе ───────────────────────────────────
\usepackage[hidelinks]{hyperref}
\usepackage{caption}

% Натписи по шаблону ФОН-а: „Табела 1: ...“ изнад табеле,
% „Слика 1. ...“ испод слике. Фонт исти као у остатку рада (12).
\captionsetup{font=normalsize,labelfont=bf,labelsep=period}
\captionsetup[table]{labelsep=colon,justification=raggedright,singlelinecheck=false}
\captionsetup[figure]{justification=centering}
\addto\captionsserbian{%
  \renewcommand{\tablename}{Табела}%
  \renewcommand{\figurename}{Слика}%
}

% ── непрекидна нумерација (1), (2), (3)... по шаблону ───────
% Шаблон приказује формуле нумерисане просто бројем у загради,
% а табеле и слике редним бројем кроз цео рад — не по поглављу.
% Зато се \numberwithin НЕ користи: подразумевана нумерација у
% класи report јесте по поглављу, па се експлицитно враћа на
% непрекидну.
\counterwithout{equation}{chapter}
\counterwithout{table}{chapter}
\counterwithout{figure}{chapter}

% ── скраћенице које се често понављају ──────────────────────
\newcommand{\RV}{\mathit{RV}}
\newcommand{\RQ}{\mathit{RQ}}
\newcommand{\BPV}{\mathit{BPV}}
\newcommand{\MedRV}{\mathit{MedRV}}
\newcommand{\TrBPV}{\mathit{TrBPV}}
\newcommand{\TPQ}{\mathit{TPQ}}
\newcommand{\RS}{\mathit{RS}}
\newcommand{\SJ}{\mathit{SJ}}
\newcommand{\QV}{\mathit{QV}}
\newcommand{\QLIKE}{\mathit{QLIKE}}
\newcommand{\RMSE}{\mathit{RMSE}}
\newcommand{\MAE}{\mathit{MAE}}
\newcommand{\VaR}{\mathit{VaR}}
\newcommand{\DM}{\mathit{DM}}
\newcommand{\BIC}{\mathit{BIC}}
\newcommand{\VIX}{\mathit{VIX}}
\newcommand{\Hit}{\mathit{Hit}}
\newcommand{\med}{\operatorname{med}}
